% ---------------------------------------------------------------------
%  CHƯƠNG 3 — PHÂN TÍCH VÀ THIẾT KẾ
% ---------------------------------------------------------------------
\chapter{Phân tích và thiết kế}
\label{ch:thietke}

\begin{muctieu}
Chương này phân tích yêu cầu bài toán, đề xuất giải pháp, thiết kế kiến trúc
hệ thống, thiết kế dữ liệu và thiết kế thành phần tri thức hoặc mô hình.
\end{muctieu}

\section{Phân tích bài toán và yêu cầu}

\begin{chuy}
\goiy{mô tả đầu vào, đầu ra của hệ thống và các yêu cầu chức năng, phi chức
năng. Có thể dùng kịch bản sử dụng \en{(use case)} để làm rõ. Hãy nêu luôn
tiêu chí thế nào là hệ thống chạy đúng.}
\end{chuy}

\begin{mau}
Hệ thống nhận đầu vào là một câu bình luận tiếng Việt và trả về nhãn cảm xúc
dự đoán kèm độ tin cậy. Yêu cầu chức năng gồm tiền xử lý văn bản, dự đoán
cảm xúc và hiển thị kết quả. Yêu cầu phi chức năng gồm thời gian phản hồi
dưới một giây cho mỗi câu và giao diện dùng được trên trình duyệt phổ thông.
\end{mau}

\section{Thiết kế kiến trúc hệ thống}

\begin{chuy}
\goiy{vẽ sơ đồ kiến trúc tổng thể gồm các khối và luồng dữ liệu. Em có thể
sửa lại sơ đồ TikZ mẫu bên dưới, hoặc chèn ảnh sơ đồ qua lệnh
\code{includegraphics}. Mọi hình đều phải được nhắc tới trong thân bài.}
\end{chuy}

Kiến trúc hệ thống được trình bày ở Hình~\ref{fig:kientruc}. Dữ liệu đi tuần
tự qua bốn khối, mỗi khối là một mô-đun độc lập trong mã nguồn.

\begin{figure}[H]
\centering
\resizebox{\textwidth}{!}{%
\begin{tikzpicture}[node distance=12mm and 14mm]
  \node[objbox] (in) {Bình luận\\ đầu vào};
  \node[clsbox, right=of in]  (pre) {Tiền xử lý\\ (tách từ, làm sạch)};
  \node[clsbox, right=of pre] (vec) {Vec-tơ hoá\\ TF-IDF};
  \node[clsbox, right=of vec] (mod) {Mô hình\\ phân lớp};
  \node[objbox, right=of mod] (out) {Nhãn cảm xúc\\ và độ tin cậy};
  \draw[flowarr] (in)  -- (pre);
  \draw[flowarr] (pre) -- (vec);
  \draw[flowarr] (vec) -- (mod);
  \draw[flowarr] (mod) -- (out);
\end{tikzpicture}}
\caption{Sơ đồ kiến trúc tổng thể của hệ thống.}
\label{fig:kientruc}
\end{figure}

\section{Thiết kế dữ liệu}

\begin{chuy}
\goiy{mô tả nguồn dữ liệu, số lượng mẫu, cấu trúc các trường và cách chia
tập huấn luyện với tập kiểm thử. Nếu có cơ sở dữ liệu thì mô tả lược đồ
bảng. Ghi rõ xuất xứ dữ liệu, đây là câu hỏi giảng viên hay hỏi khi bảo vệ.}
\end{chuy}

\begin{mau}
Bộ dữ liệu gồm 5.000 bình luận đã gán nhãn, chia theo tỉ lệ 80\% huấn luyện
và 20\% kiểm thử, chia phân tầng theo nhãn để giữ nguyên tỉ lệ các lớp. Cấu
trúc mỗi bản ghi được mô tả ở Bảng~\ref{tab:dulieu}.
\end{mau}

\begin{table}[H]
\centering
\caption{Cấu trúc bản ghi dữ liệu.}
\label{tab:dulieu}
\begin{tabular*}{\textwidth}{@{\extracolsep{\fill}}>{\bfseries\color{daunavy}}l >{\color{daumaroon}}l l}
\toprule
\rowcolor{gray100}\bfseries\color{daunavy}Trường & \bfseries\color{daumaroon}Kiểu & \bfseries\color{daunavy}Mô tả \\
\midrule
id      & integer & Khoá định danh bản ghi \\
content & text    & Nội dung bình luận \\
label   & text    & Nhãn cảm xúc: tích cực, tiêu cực, trung tính \\
\bottomrule
\end{tabular*}
\end{table}

\section{Thiết kế mô hình hoặc cơ sở tri thức}

\begin{chuy}
\goiy{đây là phần lõi của đồ án và tương ứng chuẩn đầu ra CLO3. Nếu em làm
hệ chuyên gia thì mô tả tập luật và cơ chế suy diễn. Nếu em làm học máy thì
mô tả kiến trúc mô hình, siêu tham số và cách chọn chúng. Đừng chỉ viết em
dùng thư viện nào, hãy nói rõ vì sao chọn cấu hình đó.}
\end{chuy}

\begin{mau}
Mô hình được chọn là hồi quy logistic trên đặc trưng TF-IDF với n-gram bậc
một và bậc hai, ngưỡng tần suất tối thiểu \code{min\_df} bằng 2 để loại bỏ
những từ chỉ xuất hiện một lần. Siêu tham số được chọn bằng tìm kiếm lưới
kết hợp kiểm định chéo năm phần trên tập huấn luyện, không đụng tới tập kiểm
thử trong quá trình chọn tham số.
\end{mau}

\begin{luuy}
Không được dùng tập kiểm thử để chọn siêu tham số. Làm vậy thì con số độ
chính xác báo cáo sẽ cao hơn thực tế, và em sẽ không giải thích được khi
hội đồng hỏi tới.
\end{luuy}
